\usemodule[memos]

\starttext

\title{拼音注音测试}

\section{自动注音环境测试}

\startpinyinscope
汉语拼音自动注音测试。

\color[red]{汉语拼音自动注音测试。}

{\switchtobodyfont[ss]汉语拼音自动注音测试。}

Lorem Ipsum，也称乱数假文或者哑元文本，是印
刷及排版领域所常用的虚拟文字。由于曾经一台匿名的打印机刻意打乱了一盒印刷
字体从而造出一本字体样品书，Lorem Ipsum 从西元 15 世纪起就被作为此领域的
标准文本使用。它不仅延续了五个世纪，还通过了电子排版的挑战，其雏形却依然
保存至今。在 1960 年代，『Leatraset』公司发布了印刷着 Lorem Ipsum 段落的纸
张，从而广泛普及了它的使用。最近，计算机桌面出版软件『Aldus PageMaker』
也通过同样的方式使 Lorem Ipsum 落入大众的视野。
恰恰与流行观念相反，Lorem Ipsum 并不是简简单单的随机文本。它追溯于
一篇公元前 45 年的经典拉丁著作，从而使它有着两千多年的岁数。弗吉尼亚州
Hampden­Sydney 大学拉丁系教授 Richard McClintock 曾在 Lorem Ipsum 段落中
注意到一个涵义十分隐晦的拉丁词语， 『consectetur』，通过这个单词详细查
阅跟其有关的经典文学著作原文 ，McClintock教授发掘了这个不容置疑的出
处。Lorem Ipsum 始于西塞罗(Cicero)在公元前 45 年作的『de Finibus Bonorum et
Malorum』(善恶之尽) 里 1.10.32 和 1.10.33 章节。这本书是一本关于道德理论的论
述，曾在文艺复兴时期非常流行。Lorem Ipsum 的第一行『Lorem ipsum dolor sit
amet..』节选于 1.10.32 章节。
\stoppinyinscope

\section{手动注音测试}

汉字：\xpinyin{手动注音测试}

\section{拼音输出测试}

手动注音：\xpinyin[hàn]{\bf \color[red]{汉}}

\pinyin{han4yu3pin1yin1}

\pinyin{zhong1guo2ren2min2}

\section{参数测试}

\subsection{ratio 参数（拼音字号比例）}

默认 (0.4): \xpinyin{汉}

ratio=6 (0.6): \xpinyin[ratio=6]{汉}

ratio=3 (0.3): \xpinyin[ratio=3]{汉}

ratio=8 (0.8): \xpinyin[ratio=8]{汉}

\subsection{vsep 参数（拼音汉字垂直间距）}

默认 (1em): \xpinyin{汉}

vsep=0.5em: \xpinyin[vsep=0.5em]{汉}

vsep=2em: \xpinyin[vsep=2em]{汉}

\subsection{hsep 参数（水平间距）}

hsep=10pt: \xpinyin[hsep=10pt]{汉字}

hsep=0pt: \xpinyin[hsep=0pt]{汉字}

\section{多音字测试}

\startpinyinscope
\xpinyin[chong2]{重}庆是一座美丽的城市。
\stoppinyinscope

\xpinyin{重}要

\xpinyin[chong2]{重}庆

手动设置多音字：\xpinyin[zhòng]{重}要

\subsection{multiple=yes}

\xpinyin{重庆是一座美丽的城市}

\xpinyin[multiple=yes,pysep=2pt]{重庆是一座美丽的城市}

\xpinyin[multiple=yes,pysep=8pt]{行长乐成长安}


\section{标记多音字模式}

\startpinyinscope[markpolyphone=yes]
汉语拼音自动注音测试。重庆是一座美丽的城市。
\stoppinyinscope

\xpinyin[markpolyphone=yes]{汉语拼音测试}

\setuppinyin[markpolyphone=yes]
\xpinyin{汉语拼音测试。重庆是一座美丽的城市。}

\stoptext
